\section{Discussion}

\paragraph{The silence model is the method.} The estimators are textbook.
What made them work on this array was getting ``noise'' right: which bins
contain it, which frames it can be learned from (a fixed energy gate tuned
for the processed channel missed a talker at 3\,m entirely), and when old
statistics stop describing the scene (a talker who moved between
utterances). The robot pipeline's VAD already provides the voice/silence
split, so these gains come without new sensing.

\paragraph{Distance is a room-calibrated quantity.} A 6.4\,cm array has no
geometric range resolution; the CDR measures distance relative to the
room's critical distance. This predicts, and the cross-room experiment
confirms, that a room descriptor is needed. The ReSpeaker firmware's RT$_{60}$
estimate is a candidate; its accuracy is untested. Without it, a short
calibration per room (a few labelled utterances) or coarse
near/mid/far buckets are the realistic outputs.

\paragraph{Where simulation is optimistic.} The synthetic talker is more
periodic and spectrally simpler than speech, which helps both the VAD and
PHAT; the shoebox image-source model has specular reflections and no
furniture, scattering or robot body; the diffuse field is ideal; and the
microphones are perfectly matched and positioned. Real recordings will
also reveal the mapping between raw channels, our azimuth convention and
the firmware's. We expect larger angle errors at low SNR and a weaker
within-room CDR correlation on hardware.

\paragraph{Limitations.} Small sweeps (8 clips per DOA cell, 108 distance
clips, 8 beamforming scenes) without confidence intervals; a larger
configuration is provided. One interferer, equal levels; no moving
interferer; no own-voice (TTS) handling; MVDR's noise covariance in the
beamforming experiment comes from a clean pre-onset segment.
